% TOP_PROBLEM: {"id": 30007481, "title": "Randomized k-server with squared-logarithmic ratio", "category_id": 15, "category": {"id": 15, "name": "computer_science", "display_name": "Computer Science", "description": "Computational complexity, algorithms, and theoretical CS.", "slug": "computer-science", "order_index": 15, "created_at": "2026-07-31T15:26:25.671Z"}, "status": "open"}
% FIRST_POSTED: 2026-10-06
% ATTEMPT_STATUS: solved
% OPENAI_CLAIM: {"schema_version":1,"source_repo":"https://github.com/openai/math","source_commit":"adc7f1241b42e322a6451854ab7e4b4c146bf78a","release_date":"2026-10-06","match":"full","resolution":"proved","independently_reviewed":false,"second_pass":"agreed","families":[{"family":"110","title":"Optimal-order randomized <i>k</i>-server on arbitrary metrics","subject":"Theoretical computer science","manuscripts":[{"title":"Squared-logarithmic randomized k-server on arbitrary metrics","pdf_path":"preprints/Squared-logarithmic-randomized-k-server-on-arbitrary-metrics-September-24-2026/Squared-logarithmic-randomized-k-server-on-arbitrary-metrics-September-24-2026.pdf","readme_path":"preprints/Squared-logarithmic-randomized-k-server-on-arbitrary-metrics-September-24-2026/README.md","date":"2026-09-24","theorem_ref":"preprints/Squared-logarithmic-randomized-k-server-on-arbitrary-metrics-September-24-2026/build/source/latex/sections/01_introduction.tex#thm:main"}],"lean":{"doc_path":"lean/docs/110.md","comparators":["lean/ComparatorChallenges/KServer.json","lean/ComparatorChallenges/KServer.lean","lean/ComparatorChallenges/UniformKServer.json","lean/ComparatorChallenges/UniformKServer.lean"],"declarations":[],"covers_main_theorem":true},"reasoning_trace":null}]}
\documentclass[11pt]{article}
\usepackage[margin=25mm]{geometry}
\usepackage{fontspec}
\setmainfont{Latin Modern Roman}
\usepackage{amsmath,amssymb,mathtools}
\usepackage{unicode-math}
\setmathfont{Latin Modern Math}
\usepackage{enumitem}
\usepackage{xurl,hyperref,bookmark}
\hypersetup{colorlinks=true,urlcolor=blue}
\setcounter{secnumdepth}{0}
\setlength{\parindent}{0pt}
\setlength{\parskip}{5pt}
\setlength{\emergencystretch}{3em}
\begin{document}
\section*{Randomized k-server with squared-logarithmic ratio}

\subsection{OpenAI's claim}
OpenAI presents the following claims relevant to this question, with manuscript numbers following paper order in Sources. In manuscript 1 (family 110; 2026-09-24), the asserted result is that randomized $k$-server on every metric has competitive ratio of order $\log^2(k+1)$ against oblivious requests, using one policy for every finite sequence.

\subsection{Outline of the argument}
For manuscript 1 (family 110; 2026-09-24), a hidden comparator is tracked by its posterior measure, separating changes of true metric mass from allocation costs. Protected anchors and a tree of partition labels control travel and edits, with two telescoping logarithmic budgets accounting for the final ratio.

\subsection{Scope relative to this problem}
The result supplies the conclusion in this problem with its displayed hypotheses and quantifiers. The broader statements and applications in the manuscript are kept separate from the particular question above. The supporting result is manuscript 1 (family 110; 2026-09-24).

\subsection{Formal verification}
For manuscript 1 (family 110; 2026-09-24), OpenAI lists Lean statements for this assertion: randomized $k$-server on every metric has competitive ratio of order $\log^2(k+1)$ against oblivious requests, using one policy for every finite sequence.

The selected statement includes the clause used for this problem.

The linked family Comparator checks are \texttt{KServer.lean}, \texttt{UniformKServer.lean}.

Their named declarations include \texttt{OAI.KServer.main\_theorem}, \texttt{OAI.UniformKServer.uniform\_squared\_logarithmic\_main}.

This site has not executed the Lean code. Described statement coverage is not an independent proof check.

\subsection{Status caveat}
The manuscripts were produced by an internal OpenAI model. OpenAI cautions that results without formalization can contain mistakes. This site has not independently checked the proofs; the label records the scope of OpenAI's claim.

\subsection{Sources}
\begin{itemize}
\item OpenAI, \emph{Squared-logarithmic randomized k-server on arbitrary metrics}, 2026-09-24; citation key \texttt{OAI:Squared-logarithmic-randomized-k-server-on-arbitrary-metrics-September-24-2026}. \href{https://github.com/openai/math/blob/adc7f1241b42e322a6451854ab7e4b4c146bf78a/preprints/Squared-logarithmic-randomized-k-server-on-arbitrary-metrics-September-24-2026/Squared-logarithmic-randomized-k-server-on-arbitrary-metrics-September-24-2026.pdf}{PDF (GitHub)}; \href{https://github.com/openai/math/raw/adc7f1241b42e322a6451854ab7e4b4c146bf78a/preprints/Squared-logarithmic-randomized-k-server-on-arbitrary-metrics-September-24-2026/Squared-logarithmic-randomized-k-server-on-arbitrary-metrics-September-24-2026.pdf}{PDF (direct)}; \href{https://github.com/openai/math/blob/adc7f1241b42e322a6451854ab7e4b4c146bf78a/preprints/Squared-logarithmic-randomized-k-server-on-arbitrary-metrics-September-24-2026/README.md}{Manuscript page and citation}; \href{https://github.com/openai/math/blob/main/preprints/Squared-logarithmic-randomized-k-server-on-arbitrary-metrics-September-24-2026/Squared-logarithmic-randomized-k-server-on-arbitrary-metrics-September-24-2026.pdf}{Latest version}.
\item Family 110: \href{https://github.com/openai/math/blob/adc7f1241b42e322a6451854ab7e4b4c146bf78a/lean/docs/110.md}{Lean scope}.
\item \href{https://github.com/openai/math/blob/adc7f1241b42e322a6451854ab7e4b4c146bf78a/lean/ComparatorChallenges/KServer.json}{Comparator KServer.json}.
\item \href{https://github.com/openai/math/blob/adc7f1241b42e322a6451854ab7e4b4c146bf78a/lean/ComparatorChallenges/KServer.lean}{Comparator KServer.lean}.
\item \href{https://github.com/openai/math/blob/adc7f1241b42e322a6451854ab7e4b4c146bf78a/lean/ComparatorChallenges/UniformKServer.json}{Comparator UniformKServer.json}.
\item \href{https://github.com/openai/math/blob/adc7f1241b42e322a6451854ab7e4b4c146bf78a/lean/ComparatorChallenges/UniformKServer.lean}{Comparator UniformKServer.lean}.
\item \href{https://github.com/openai/math/blob/adc7f1241b42e322a6451854ab7e4b4c146bf78a/CONTENTS.md}{OpenAI catalogue entry}.
\end{itemize}
\end{document}
